Depth reduction turns an oblivious tape computation of matrix transposition into an index code whose side information consists of bounded dependency cones. For fixed-depth cones of linear-time native schedules, we prove that, after paid degree trimming, only $O(L/\log d)$ matched input and demand ports lie within any fixed incidence distance. Here $L=d^2b$ is the bit length of a $d\times d$ matrix. The theorem uses the full syntactic dependency graph and does not assume that the computation is correct.

We obtain half-rate bounds for bounded complex linear recovery, for arbitrary finite-alphabet codes on uniformly hyperfinite incidence graphs, and for directed cyclic components with at most two internal inputs. Cheap native executions show that neither hyperfiniteness nor sublinear reduction to those cyclic components is universal. Other native examples retain many connected demand pairs while admitting an all-code deficit independent of fixed cone depth; a reusable row-word theorem isolates the interface behind that deficit.

An algebraic appendix gives matrix-label certificates uniform in vector width and applies them to an expanding characteristic-two family. A central-layer argument transfers this bound to nonlinear codes with finite-$2$-group endomorphic recovery expressions, including a quantitative payment for exceptional rows.

Finite counting and algebraic interfaces are verified in Lean. The native bridge and its applications are partly formalized; the analytic, cyclic-component, group-transfer, and graph-construction arguments have written proofs with their status stated individually. The general finite-alphabet coding question, unrestricted cheap depth reduction, and a paid treatment of adaptive schedules remain open. No unrestricted transposition time bound follows here.
